\documentclass[11pt]{article}

\usepackage[margin=1in]{geometry}
\usepackage{amsmath,amssymb,amsthm}
\usepackage{graphicx}
\usepackage{hyperref}
\usepackage{xurl}
\usepackage{microtype}
\usepackage{enumitem}
\usepackage{booktabs}
\usepackage{listings}
\usepackage[nameinlink,noabbrev]{cleveref}

\graphicspath{{../figures/}}

\newtheorem{theorem}{Theorem}
\newtheorem{lemma}{Lemma}
\newtheorem{proposition}{Proposition}
\newtheorem{corollary}{Corollary}
\theoremstyle{definition}
\newtheorem{definition}{Definition}
\theoremstyle{remark}
\newtheorem{remark}{Remark}

\newcommand{\TV}{\mathrm{TV}}
\newcommand{\eps}{\varepsilon}

\title{PSC-Q (Addendum):\\Supplementary Evidence and Operator Microstudies}
\author{}
\date{February 25, 2026}

\begin{document}
\maketitle

\begin{abstract}
This companion note preserves the PSC-Q microstudies and operator-facing ``gotchas'' that were previously shipped as a split supplement.
The main PSC-Q paper focuses on the mechanism interface and declarations; this addendum collects short empirical illustrations and checklists.
Figures are omitted in this source-only release; corresponding plots and raw tables are available on request.\cite{series:artifactpolicy}
\end{abstract}

\paragraph{Scope.}
Addendum to PSC-Q.\cite{series:pscq}

\paragraph{Imports.}
We import Congestion-EQ as the target contract and PSC-Q as the mechanism surface.\cite{series:congeq,series:pscq}

\paragraph{Exports.}
A compact set of sanity checks and microstudy patterns for implementers.

\section{Supplementary evidence and microstudies (formerly a split supplement)}\label{app:microstudies}
This appendix preserves the operator-facing microstudies and ``gotchas'' that were previously shipped as a separate
PSC-Q supplement in the split archive. The unified archive keeps a single-file rule, so we inline the salient parts here.
Figures are omitted in this source-only release; the corresponding plots and raw tables are available on request.

\subsection{Additive padding can worsen congestion-TV in the high-utilization knee}
The main text emphasizes \emph{substitution} (idle-filling) because it can reduce congestion witnesses without increasing mean load.
In contrast, \emph{additive} padding can increase utilization and move the system closer to the slack ``knee,''
at which point small secret-correlated load differences yield large changes in waiting-time and RTT laws.
Fig.~\ref{fig:worsens} (from the research archive) illustrates a simple regime in which additive padding increases the measured congestion-TV.

\begin{figure}[t]
\centering
\fbox{\begin{minipage}{0.92\linewidth}\centering\small Figure omitted in source-only release. Requested file: \texttt{fig\_padding\_worsens\_tv.pdf}.\end{minipage}}
\caption{Archive microstudy: a ``padding can worsen'' regime when additive padding pushes utilization upward.}
\label{fig:worsens}
\end{figure}

\subsection{Partial idle-filling is often enough}
Even partial substitution (filling only a fraction of idle slots) can materially reduce congestion-TV without increasing mean load.
Fig.~\ref{fig:idle} summarizes an archive sweep supporting this empirical claim.

\begin{figure}[t]
\centering
\fbox{\begin{minipage}{0.92\linewidth}\centering\small Figure omitted in source-only release. Requested file: \texttt{fig\_idle\_padding\_partial\_tv.pdf}.\end{minipage}}
\caption{Archive microstudy: partial idle-filling can already reduce congestion-TV without increasing mean load.}
\label{fig:idle}
\end{figure}

\subsection{Substitution must avoid service bias (engineering pitfall)}
Substitution is only ``free'' if cover is confined to otherwise-idle service opportunities.
If cover work competes with real work (shared queue or shared downstream resource), it can bias service and increase delay.
Fig.~\ref{fig:servicebias} illustrates this pitfall under a simple mixed-queue implementation.

\begin{figure}[t]
\centering
\fbox{\begin{minipage}{0.92\linewidth}\centering\small Figure omitted in source-only release. Requested file: \texttt{fig\_service\_bias\_substitution.pdf}.\end{minipage}}
\caption{Archive microstudy: a service-bias pitfall when ``cover'' is implemented as competing work rather than strict idle-filling.
Implementation takeaway: implement substitution as \emph{preemptible} idle-filling (or as a separate low-priority queue that never delays real work).}
\label{fig:servicebias}
\end{figure}

\paragraph{Implementation checklist (non-exhaustive).}
\begin{itemize}[leftmargin=1.2em]
\item Define at most one service opportunity per slot; do not create extra opportunities when real work is present.
\item Treat cover as cancelable/preemptible: if real work arrives, cancel or defer cover so it never delays real work.
\item If a separate queue is used, enforce strict priority for real work and bound head-of-line blocking (e.g., fixed-size cover frames).
\item Evaluate congestion witnesses across hops (RTT time series, queue delay, loss/ECN) to avoid bottleneck shift.
\end{itemize}

\subsection{Spectral surrogate evidence (why a ``spectral tier'' exists)}
The PSC-Q main text treats a TV-tier contract and flags periodic/spectral tests as a realistic additional tier.
Fig.~\ref{fig:spectral} is an archive ``spectral surrogate'' plot illustrating that declared dither attenuates periodic witnesses,
supporting the motivation for a two-tier view.

\begin{figure}[t]
\centering
\fbox{\begin{minipage}{0.92\linewidth}\centering\small Figure omitted in source-only release. Requested file: \texttt{fig\_spectral\_dither\_attenuation.pdf}.\end{minipage}}
\caption{Archive evidence: dither attenuates periodic witnesses in a spectral surrogate.}
\label{fig:spectral}
\end{figure}

\subsection{Interval-robust compilation against a utilization range}
The ``range compilation'' note in the main text can be made more explicit:
compile $(\delta,\alpha)$ against a declared utilization interval $\rho\in[\rho_{\min},\rho_{\max}]$ (rather than a point estimate),
and size for $\rho_{\max}$ so the deployment remains within budget whenever the monitored regime stays within the published interval.
Fig.~\ref{fig:utilrangeopt} illustrates an archive sweep in which the chosen $(\delta,\alpha)$ tightens as $\rho_{\max}$ increases.

\begin{figure}[t]
\centering
\fbox{\begin{minipage}{0.92\linewidth}\centering\small Figure omitted in source-only release. Requested file: \texttt{fig\_utilization\_range\_optimizer.pdf}.\end{minipage}}
\caption{Archive evidence: interval-robust compilation tightens as the assumed upper utilization bound increases.}
\label{fig:utilrangeopt}
\end{figure}

\begin{thebibliography}{99}

\bibitem{series:pscq}
Anonymous.
\newblock PSC-Q: Public Service Calendars and Slack-Harvesting Substitution Padding for Congestion-Private Anonymous Overlays.
\newblock Congestion series Paper~2, 2026. (This archive.)

\bibitem{series:congeq}
Anonymous.
\newblock Congestion-EQ: A Congestion Witness Contract for Anonymous Overlays.
\newblock Congestion series Paper~1, 2026. (This archive.)

\bibitem{series:artifactpolicy}
Anonymous.
\newblock Artifact policy and supporting materials for the anonymity/AnonDHT series.
\newblock 2026.
\end{thebibliography}

\end{document}
